% ---------------------------------------------------------------------------
% RQ5 results section -- FaceCon FYP (F26-082).  % ---------------------------------------------------------------------------
% RQ5 results section -- FaceCon FYP (F26-082).  % ---------------------------------------------------------------------------
% RQ5 results section -- FaceCon FYP (F26-082).  % ---------------------------------------------------------------------------
% RQ5 results section -- FaceCon FYP (F26-082).  \input{rq5_results} or paste.
% Source of every number: gpu/report_back.json and gpu/controls.json
% (RTX A4000 run of `python -m gpu.run_all` and `python -m gpu.controls --power`).
% ALL DATA ARE SIMULATED. Nothing here is a finding about real faces.
% Replaces the earlier draft, whose conclusions were a feature-pipeline artefact.
% ---------------------------------------------------------------------------

\subsection{RQ5: Do more complex models beat simple baselines?}
\label{sec:rq5}

\paragraph{Scope.}
All results in this section come from simulated data in which a known share of
facial-shape variance (20\,\% genetic, 35\,\% demographic, 45\,\% noise) was
planted. They test whether the pipeline and models behave correctly; they say
nothing about how well genotype predicts real faces.

\paragraph{Protocol.}
The matched cohort has $N=3{,}118$ subjects (the size of the real analysis
sample), a 5{,}000-SNP panel and a 50-dimensional shape target, split
$2{,}182/467/469$ at subject level. All scaling, feature selection and model
fitting use the training split only; each deep model is scored at its
best-validation checkpoint. We report mean test $R^2$ over the 50 shape
dimensions. The $\pm$ values are standard deviations over three training seeds
on one dataset and one split, so they capture initialisation noise only.
Bootstrap intervals on this test set are roughly $\pm0.01$, and differences
smaller than that should not be read as real.

\paragraph{Feature representation mattered more than architecture.}
Our first experiments gave every model 256 SNP principal components. A control
with ground-truth causal SNPs showed that this step keeps only $+0.004$ of the
$+0.173$ $R^2$ that the SNPs can add beyond covariates and ancestry components.
We therefore added a second representation: the 256 SNPs most associated with
shape on the training split after removing demographic effects. Table
\ref{tab:rq5} reports both.

\begin{table}[t]
\centering
\caption{Matched synthetic cohort ($N=3{,}118$): test $R^2$, mean $\pm$ s.d.\
over three training seeds. The demographic baseline is $0.317$ in both columns.}
\label{tab:rq5}
\begin{tabular}{lrr}
\toprule
Model (with demographics) & SNP principal components & Screened SNPs \\
\midrule
Demographics only (baseline) & $0.317$ & $0.317$ \\
MLP                          & $0.274\pm0.001$ & $\mathbf{0.396\pm0.002}$ \\
Contrastive encoder          & $0.220\pm0.004$ & $0.364\pm0.001$ \\
Transformer SNP encoder      & $0.311\pm0.005$ & $0.319\pm0.002$ \\
Conditional diffusion        & $-0.368\pm0.006$ & $0.081\pm0.006$ \\
\midrule
\multicolumn{2}{l}{Linear reference: covariates + ancestry components} & $0.315\pm0.003$ \\
\multicolumn{2}{l}{Linear reference: + screened SNPs, ridge (tuned)}   & $0.488\pm0.004$ \\
\multicolumn{2}{l}{Oracle: + true causal SNPs}                         & $0.488\pm0.005$ \\
\bottomrule
\end{tabular}
\end{table}

With principal components no deep model exceeds the demographic baseline. With
screened SNPs the MLP ($0.396$) and the contrastive encoder ($0.364$) both
exceed it clearly, and mean mesh error for the MLP falls from $3.00$\,mm
(demographics) to $2.79$\,mm. The Transformer stays at the baseline ($0.319$)
and the diffusion decoder remains far below it.

\paragraph{Answer to RQ5.}
On this simulated task, complexity did not pay. The best deep model recovers
$+0.079$ over demographics, less than half of the $+0.173$ available, while a
tuned ridge regression on screened SNPs reaches $0.488$ and matches the oracle.
The linear reference rows come from a separate script with five dataset seeds
and a tuned number of SNPs, so the comparison with the deep models is indicative
rather than exact; the gap ($0.09$) is nevertheless far larger than any
plausible protocol difference. The planted effects are linear and additive, so
a linear model is the correctly specified one here. This result therefore shows
that the deep models are less data-efficient on a linear signal, not that deep
models are worse in general.

\subsection{Controls}
\label{sec:rq5-controls}

\begin{table}[t]
\centering
\caption{Controls on the matched cohort. Positive control: planted signal must
be found. Negative controls: no signal must be reported when none exists.}
\label{tab:rq5-controls}
\begin{tabular}{lr}
\toprule
Condition & $\Delta R^2$ from adding SNPs \\
\midrule
Screened SNPs, planted share 0.20 (5 seeds)   & $+0.173$ (per-seed 95\,\% CIs within $[+0.159,+0.188]$) \\
Genotypes permuted across subjects (5 seeds)  & $-0.001$ \\
Planted genetic share set to zero (3 seeds)   & $-0.001$ \\
\bottomrule
\end{tabular}
\end{table}

The pipeline finds the signal when it is planted and reports none when the
genotype--face pairing is broken or the genetic share is zero (Table
\ref{tab:rq5-controls}). The permutation control was also run on the deep
models: the Transformer scored $0.316$, the MLP $0.252$ and the contrastive
encoder $0.247$, against $0.317$ for demographics. Two points follow. First, no
model gains from uninformative SNPs, which rules out leakage. Second, the MLP
and contrastive encoder lose about $0.07$ when given uninformative SNPs, so on
data with weak genetic signal they could fall below the demographic baseline.
The Transformer's permuted score equals its unpermuted score, which indicates
that at this sample size it is not using the SNPs at all.

\subsection{Effect of sample size}
\label{sec:rq5-scaling}

To ask whether more subjects would help, we generated a separate dataset of
$N=10{,}000$ with the same settings and trained on nested subsets, keeping one
fixed test set of $1{,}500$ (Table \ref{tab:rq5-scaling}). This is a scaling
experiment only: the real cohort has $3{,}118$ subjects, so results beyond the
first row describe a cohort we will not have.

\begin{table}[t]
\centering
\caption{Learning curve on the $N=10{,}000$ scaling dataset, screened SNPs,
test $R^2$ (mean $\pm$ s.d.\ over three training seeds).}
\label{tab:rq5-scaling}
\begin{tabular}{rrrrrr}
\toprule
Train $N$ & Demographics & MLP & Contrastive & Transformer & Diffusion \\
\midrule
$2{,}182$ & $0.312$ & $0.389\pm0.001$ & $0.349\pm0.002$ & $0.310\pm0.001$ & $0.058\pm0.003$ \\
$4{,}000$ & $0.313$ & $0.417\pm0.001$ & $0.392\pm0.001$ & $0.316\pm0.002$ & $0.049\pm0.012$ \\
$7{,}000$ & $0.314$ & $0.435\pm0.001$ & $0.428\pm0.000$ & $0.338\pm0.003$ & $0.028\pm0.011$ \\
\bottomrule
\end{tabular}
\end{table}

The demographic baseline is flat, as expected for a five-parameter linear model.
The MLP gains $+0.046$ and the contrastive encoder $+0.079$ from roughly three
times the data, and the Transformer begins to exceed the baseline only at
$7{,}000$. The diffusion decoder does not improve; we have not diagnosed why and
report it as an open negative result. A companion power grid with the linear
screen gives the same picture: at $N=3{,}118$ the SNPs add $+0.070$ when the
planted genetic share is 0.10 and $+0.177$ when it is 0.20, and nothing when it
is zero.

\paragraph{Limits.}
The simulation is easier than real data in ways that favour detection: 5{,}000
unlinked SNPs instead of roughly $10^6$ correlated ones, 150 causal SNPs that
each affect every shape dimension, and three cleanly separated ancestry groups.
The power figures are therefore an optimistic upper bound. What this section
establishes is that the pipeline is correct, leak-free and sensitive to planted
signal, and that the model comparison is ready to be repeated unchanged on real
data.
 or paste.
% Source of every number: gpu/report_back.json and gpu/controls.json
% (RTX A4000 run of `python -m gpu.run_all` and `python -m gpu.controls --power`).
% ALL DATA ARE SIMULATED. Nothing here is a finding about real faces.
% Replaces the earlier draft, whose conclusions were a feature-pipeline artefact.
% ---------------------------------------------------------------------------

\subsection{RQ5: Do more complex models beat simple baselines?}
\label{sec:rq5}

\paragraph{Scope.}
All results in this section come from simulated data in which a known share of
facial-shape variance (20\,\% genetic, 35\,\% demographic, 45\,\% noise) was
planted. They test whether the pipeline and models behave correctly; they say
nothing about how well genotype predicts real faces.

\paragraph{Protocol.}
The matched cohort has $N=3{,}118$ subjects (the size of the real analysis
sample), a 5{,}000-SNP panel and a 50-dimensional shape target, split
$2{,}182/467/469$ at subject level. All scaling, feature selection and model
fitting use the training split only; each deep model is scored at its
best-validation checkpoint. We report mean test $R^2$ over the 50 shape
dimensions. The $\pm$ values are standard deviations over three training seeds
on one dataset and one split, so they capture initialisation noise only.
Bootstrap intervals on this test set are roughly $\pm0.01$, and differences
smaller than that should not be read as real.

\paragraph{Feature representation mattered more than architecture.}
Our first experiments gave every model 256 SNP principal components. A control
with ground-truth causal SNPs showed that this step keeps only $+0.004$ of the
$+0.173$ $R^2$ that the SNPs can add beyond covariates and ancestry components.
We therefore added a second representation: the 256 SNPs most associated with
shape on the training split after removing demographic effects. Table
\ref{tab:rq5} reports both.

\begin{table}[t]
\centering
\caption{Matched synthetic cohort ($N=3{,}118$): test $R^2$, mean $\pm$ s.d.\
over three training seeds. The demographic baseline is $0.317$ in both columns.}
\label{tab:rq5}
\begin{tabular}{lrr}
\toprule
Model (with demographics) & SNP principal components & Screened SNPs \\
\midrule
Demographics only (baseline) & $0.317$ & $0.317$ \\
MLP                          & $0.274\pm0.001$ & $\mathbf{0.396\pm0.002}$ \\
Contrastive encoder          & $0.220\pm0.004$ & $0.364\pm0.001$ \\
Transformer SNP encoder      & $0.311\pm0.005$ & $0.319\pm0.002$ \\
Conditional diffusion        & $-0.368\pm0.006$ & $0.081\pm0.006$ \\
\midrule
\multicolumn{2}{l}{Linear reference: covariates + ancestry components} & $0.315\pm0.003$ \\
\multicolumn{2}{l}{Linear reference: + screened SNPs, ridge (tuned)}   & $0.488\pm0.004$ \\
\multicolumn{2}{l}{Oracle: + true causal SNPs}                         & $0.488\pm0.005$ \\
\bottomrule
\end{tabular}
\end{table}

With principal components no deep model exceeds the demographic baseline. With
screened SNPs the MLP ($0.396$) and the contrastive encoder ($0.364$) both
exceed it clearly, and mean mesh error for the MLP falls from $3.00$\,mm
(demographics) to $2.79$\,mm. The Transformer stays at the baseline ($0.319$)
and the diffusion decoder remains far below it.

\paragraph{Answer to RQ5.}
On this simulated task, complexity did not pay. The best deep model recovers
$+0.079$ over demographics, less than half of the $+0.173$ available, while a
tuned ridge regression on screened SNPs reaches $0.488$ and matches the oracle.
The linear reference rows come from a separate script with five dataset seeds
and a tuned number of SNPs, so the comparison with the deep models is indicative
rather than exact; the gap ($0.09$) is nevertheless far larger than any
plausible protocol difference. The planted effects are linear and additive, so
a linear model is the correctly specified one here. This result therefore shows
that the deep models are less data-efficient on a linear signal, not that deep
models are worse in general.

\subsection{Controls}
\label{sec:rq5-controls}

\begin{table}[t]
\centering
\caption{Controls on the matched cohort. Positive control: planted signal must
be found. Negative controls: no signal must be reported when none exists.}
\label{tab:rq5-controls}
\begin{tabular}{lr}
\toprule
Condition & $\Delta R^2$ from adding SNPs \\
\midrule
Screened SNPs, planted share 0.20 (5 seeds)   & $+0.173$ (per-seed 95\,\% CIs within $[+0.159,+0.188]$) \\
Genotypes permuted across subjects (5 seeds)  & $-0.001$ \\
Planted genetic share set to zero (3 seeds)   & $-0.001$ \\
\bottomrule
\end{tabular}
\end{table}

The pipeline finds the signal when it is planted and reports none when the
genotype--face pairing is broken or the genetic share is zero (Table
\ref{tab:rq5-controls}). The permutation control was also run on the deep
models: the Transformer scored $0.316$, the MLP $0.252$ and the contrastive
encoder $0.247$, against $0.317$ for demographics. Two points follow. First, no
model gains from uninformative SNPs, which rules out leakage. Second, the MLP
and contrastive encoder lose about $0.07$ when given uninformative SNPs, so on
data with weak genetic signal they could fall below the demographic baseline.
The Transformer's permuted score equals its unpermuted score, which indicates
that at this sample size it is not using the SNPs at all.

\subsection{Effect of sample size}
\label{sec:rq5-scaling}

To ask whether more subjects would help, we generated a separate dataset of
$N=10{,}000$ with the same settings and trained on nested subsets, keeping one
fixed test set of $1{,}500$ (Table \ref{tab:rq5-scaling}). This is a scaling
experiment only: the real cohort has $3{,}118$ subjects, so results beyond the
first row describe a cohort we will not have.

\begin{table}[t]
\centering
\caption{Learning curve on the $N=10{,}000$ scaling dataset, screened SNPs,
test $R^2$ (mean $\pm$ s.d.\ over three training seeds).}
\label{tab:rq5-scaling}
\begin{tabular}{rrrrrr}
\toprule
Train $N$ & Demographics & MLP & Contrastive & Transformer & Diffusion \\
\midrule
$2{,}182$ & $0.312$ & $0.389\pm0.001$ & $0.349\pm0.002$ & $0.310\pm0.001$ & $0.058\pm0.003$ \\
$4{,}000$ & $0.313$ & $0.417\pm0.001$ & $0.392\pm0.001$ & $0.316\pm0.002$ & $0.049\pm0.012$ \\
$7{,}000$ & $0.314$ & $0.435\pm0.001$ & $0.428\pm0.000$ & $0.338\pm0.003$ & $0.028\pm0.011$ \\
\bottomrule
\end{tabular}
\end{table}

The demographic baseline is flat, as expected for a five-parameter linear model.
The MLP gains $+0.046$ and the contrastive encoder $+0.079$ from roughly three
times the data, and the Transformer begins to exceed the baseline only at
$7{,}000$. The diffusion decoder does not improve; we have not diagnosed why and
report it as an open negative result. A companion power grid with the linear
screen gives the same picture: at $N=3{,}118$ the SNPs add $+0.070$ when the
planted genetic share is 0.10 and $+0.177$ when it is 0.20, and nothing when it
is zero.

\paragraph{Limits.}
The simulation is easier than real data in ways that favour detection: 5{,}000
unlinked SNPs instead of roughly $10^6$ correlated ones, 150 causal SNPs that
each affect every shape dimension, and three cleanly separated ancestry groups.
The power figures are therefore an optimistic upper bound. What this section
establishes is that the pipeline is correct, leak-free and sensitive to planted
signal, and that the model comparison is ready to be repeated unchanged on real
data.
 or paste.
% Source of every number: gpu/report_back.json and gpu/controls.json
% (RTX A4000 run of `python -m gpu.run_all` and `python -m gpu.controls --power`).
% ALL DATA ARE SIMULATED. Nothing here is a finding about real faces.
% Replaces the earlier draft, whose conclusions were a feature-pipeline artefact.
% ---------------------------------------------------------------------------

\subsection{RQ5: Do more complex models beat simple baselines?}
\label{sec:rq5}

\paragraph{Scope.}
All results in this section come from simulated data in which a known share of
facial-shape variance (20\,\% genetic, 35\,\% demographic, 45\,\% noise) was
planted. They test whether the pipeline and models behave correctly; they say
nothing about how well genotype predicts real faces.

\paragraph{Protocol.}
The matched cohort has $N=3{,}118$ subjects (the size of the real analysis
sample), a 5{,}000-SNP panel and a 50-dimensional shape target, split
$2{,}182/467/469$ at subject level. All scaling, feature selection and model
fitting use the training split only; each deep model is scored at its
best-validation checkpoint. We report mean test $R^2$ over the 50 shape
dimensions. The $\pm$ values are standard deviations over three training seeds
on one dataset and one split, so they capture initialisation noise only.
Bootstrap intervals on this test set are roughly $\pm0.01$, and differences
smaller than that should not be read as real.

\paragraph{Feature representation mattered more than architecture.}
Our first experiments gave every model 256 SNP principal components. A control
with ground-truth causal SNPs showed that this step keeps only $+0.004$ of the
$+0.173$ $R^2$ that the SNPs can add beyond covariates and ancestry components.
We therefore added a second representation: the 256 SNPs most associated with
shape on the training split after removing demographic effects. Table
\ref{tab:rq5} reports both.

\begin{table}[t]
\centering
\caption{Matched synthetic cohort ($N=3{,}118$): test $R^2$, mean $\pm$ s.d.\
over three training seeds. The demographic baseline is $0.317$ in both columns.}
\label{tab:rq5}
\begin{tabular}{lrr}
\toprule
Model (with demographics) & SNP principal components & Screened SNPs \\
\midrule
Demographics only (baseline) & $0.317$ & $0.317$ \\
MLP                          & $0.274\pm0.001$ & $\mathbf{0.396\pm0.002}$ \\
Contrastive encoder          & $0.220\pm0.004$ & $0.364\pm0.001$ \\
Transformer SNP encoder      & $0.311\pm0.005$ & $0.319\pm0.002$ \\
Conditional diffusion        & $-0.368\pm0.006$ & $0.081\pm0.006$ \\
\midrule
\multicolumn{2}{l}{Linear reference: covariates + ancestry components} & $0.315\pm0.003$ \\
\multicolumn{2}{l}{Linear reference: + screened SNPs, ridge (tuned)}   & $0.488\pm0.004$ \\
\multicolumn{2}{l}{Oracle: + true causal SNPs}                         & $0.488\pm0.005$ \\
\bottomrule
\end{tabular}
\end{table}

With principal components no deep model exceeds the demographic baseline. With
screened SNPs the MLP ($0.396$) and the contrastive encoder ($0.364$) both
exceed it clearly, and mean mesh error for the MLP falls from $3.00$\,mm
(demographics) to $2.79$\,mm. The Transformer stays at the baseline ($0.319$)
and the diffusion decoder remains far below it.

\paragraph{Answer to RQ5.}
On this simulated task, complexity did not pay. The best deep model recovers
$+0.079$ over demographics, less than half of the $+0.173$ available, while a
tuned ridge regression on screened SNPs reaches $0.488$ and matches the oracle.
The linear reference rows come from a separate script with five dataset seeds
and a tuned number of SNPs, so the comparison with the deep models is indicative
rather than exact; the gap ($0.09$) is nevertheless far larger than any
plausible protocol difference. The planted effects are linear and additive, so
a linear model is the correctly specified one here. This result therefore shows
that the deep models are less data-efficient on a linear signal, not that deep
models are worse in general.

\subsection{Controls}
\label{sec:rq5-controls}

\begin{table}[t]
\centering
\caption{Controls on the matched cohort. Positive control: planted signal must
be found. Negative controls: no signal must be reported when none exists.}
\label{tab:rq5-controls}
\begin{tabular}{lr}
\toprule
Condition & $\Delta R^2$ from adding SNPs \\
\midrule
Screened SNPs, planted share 0.20 (5 seeds)   & $+0.173$ (per-seed 95\,\% CIs within $[+0.159,+0.188]$) \\
Genotypes permuted across subjects (5 seeds)  & $-0.001$ \\
Planted genetic share set to zero (3 seeds)   & $-0.001$ \\
\bottomrule
\end{tabular}
\end{table}

The pipeline finds the signal when it is planted and reports none when the
genotype--face pairing is broken or the genetic share is zero (Table
\ref{tab:rq5-controls}). The permutation control was also run on the deep
models: the Transformer scored $0.316$, the MLP $0.252$ and the contrastive
encoder $0.247$, against $0.317$ for demographics. Two points follow. First, no
model gains from uninformative SNPs, which rules out leakage. Second, the MLP
and contrastive encoder lose about $0.07$ when given uninformative SNPs, so on
data with weak genetic signal they could fall below the demographic baseline.
The Transformer's permuted score equals its unpermuted score, which indicates
that at this sample size it is not using the SNPs at all.

\subsection{Effect of sample size}
\label{sec:rq5-scaling}

To ask whether more subjects would help, we generated a separate dataset of
$N=10{,}000$ with the same settings and trained on nested subsets, keeping one
fixed test set of $1{,}500$ (Table \ref{tab:rq5-scaling}). This is a scaling
experiment only: the real cohort has $3{,}118$ subjects, so results beyond the
first row describe a cohort we will not have.

\begin{table}[t]
\centering
\caption{Learning curve on the $N=10{,}000$ scaling dataset, screened SNPs,
test $R^2$ (mean $\pm$ s.d.\ over three training seeds).}
\label{tab:rq5-scaling}
\begin{tabular}{rrrrrr}
\toprule
Train $N$ & Demographics & MLP & Contrastive & Transformer & Diffusion \\
\midrule
$2{,}182$ & $0.312$ & $0.389\pm0.001$ & $0.349\pm0.002$ & $0.310\pm0.001$ & $0.058\pm0.003$ \\
$4{,}000$ & $0.313$ & $0.417\pm0.001$ & $0.392\pm0.001$ & $0.316\pm0.002$ & $0.049\pm0.012$ \\
$7{,}000$ & $0.314$ & $0.435\pm0.001$ & $0.428\pm0.000$ & $0.338\pm0.003$ & $0.028\pm0.011$ \\
\bottomrule
\end{tabular}
\end{table}

The demographic baseline is flat, as expected for a five-parameter linear model.
The MLP gains $+0.046$ and the contrastive encoder $+0.079$ from roughly three
times the data, and the Transformer begins to exceed the baseline only at
$7{,}000$. The diffusion decoder does not improve; we have not diagnosed why and
report it as an open negative result. A companion power grid with the linear
screen gives the same picture: at $N=3{,}118$ the SNPs add $+0.070$ when the
planted genetic share is 0.10 and $+0.177$ when it is 0.20, and nothing when it
is zero.

\paragraph{Limits.}
The simulation is easier than real data in ways that favour detection: 5{,}000
unlinked SNPs instead of roughly $10^6$ correlated ones, 150 causal SNPs that
each affect every shape dimension, and three cleanly separated ancestry groups.
The power figures are therefore an optimistic upper bound. What this section
establishes is that the pipeline is correct, leak-free and sensitive to planted
signal, and that the model comparison is ready to be repeated unchanged on real
data.
 or paste.
% Source of every number: gpu/report_back.json and gpu/controls.json
% (RTX A4000 run of `python -m gpu.run_all` and `python -m gpu.controls --power`).
% ALL DATA ARE SIMULATED. Nothing here is a finding about real faces.
% Replaces the earlier draft, whose conclusions were a feature-pipeline artefact.
% ---------------------------------------------------------------------------

\subsection{RQ5: Do more complex models beat simple baselines?}
\label{sec:rq5}

\paragraph{Scope.}
All results in this section come from simulated data in which a known share of
facial-shape variance (20\,\% genetic, 35\,\% demographic, 45\,\% noise) was
planted. They test whether the pipeline and models behave correctly; they say
nothing about how well genotype predicts real faces.

\paragraph{Protocol.}
The matched cohort has $N=3{,}118$ subjects (the size of the real analysis
sample), a 5{,}000-SNP panel and a 50-dimensional shape target, split
$2{,}182/467/469$ at subject level. All scaling, feature selection and model
fitting use the training split only; each deep model is scored at its
best-validation checkpoint. We report mean test $R^2$ over the 50 shape
dimensions. The $\pm$ values are standard deviations over three training seeds
on one dataset and one split, so they capture initialisation noise only.
Bootstrap intervals on this test set are roughly $\pm0.01$, and differences
smaller than that should not be read as real.

\paragraph{Feature representation mattered more than architecture.}
Our first experiments gave every model 256 SNP principal components. A control
with ground-truth causal SNPs showed that this step keeps only $+0.004$ of the
$+0.173$ $R^2$ that the SNPs can add beyond covariates and ancestry components.
We therefore added a second representation: the 256 SNPs most associated with
shape on the training split after removing demographic effects. Table
\ref{tab:rq5} reports both.

\begin{table}[t]
\centering
\caption{Matched synthetic cohort ($N=3{,}118$): test $R^2$, mean $\pm$ s.d.\
over three training seeds. The demographic baseline is $0.317$ in both columns.}
\label{tab:rq5}
\begin{tabular}{lrr}
\toprule
Model (with demographics) & SNP principal components & Screened SNPs \\
\midrule
Demographics only (baseline) & $0.317$ & $0.317$ \\
MLP                          & $0.274\pm0.001$ & $\mathbf{0.396\pm0.002}$ \\
Contrastive encoder          & $0.220\pm0.004$ & $0.364\pm0.001$ \\
Transformer SNP encoder      & $0.311\pm0.005$ & $0.319\pm0.002$ \\
Conditional diffusion        & $-0.368\pm0.006$ & $0.081\pm0.006$ \\
\midrule
\multicolumn{2}{l}{Linear reference: covariates + ancestry components} & $0.315\pm0.003$ \\
\multicolumn{2}{l}{Linear reference: + screened SNPs, ridge (tuned)}   & $0.488\pm0.004$ \\
\multicolumn{2}{l}{Oracle: + true causal SNPs}                         & $0.488\pm0.005$ \\
\bottomrule
\end{tabular}
\end{table}

With principal components no deep model exceeds the demographic baseline. With
screened SNPs the MLP ($0.396$) and the contrastive encoder ($0.364$) both
exceed it clearly, and mean mesh error for the MLP falls from $3.00$\,mm
(demographics) to $2.79$\,mm. The Transformer stays at the baseline ($0.319$)
and the diffusion decoder remains far below it.

\paragraph{Answer to RQ5.}
On this simulated task, complexity did not pay. The best deep model recovers
$+0.079$ over demographics, less than half of the $+0.173$ available, while a
tuned ridge regression on screened SNPs reaches $0.488$ and matches the oracle.
The linear reference rows come from a separate script with five dataset seeds
and a tuned number of SNPs, so the comparison with the deep models is indicative
rather than exact; the gap ($0.09$) is nevertheless far larger than any
plausible protocol difference. The planted effects are linear and additive, so
a linear model is the correctly specified one here. This result therefore shows
that the deep models are less data-efficient on a linear signal, not that deep
models are worse in general.

\subsection{Controls}
\label{sec:rq5-controls}

\begin{table}[t]
\centering
\caption{Controls on the matched cohort. Positive control: planted signal must
be found. Negative controls: no signal must be reported when none exists.}
\label{tab:rq5-controls}
\begin{tabular}{lr}
\toprule
Condition & $\Delta R^2$ from adding SNPs \\
\midrule
Screened SNPs, planted share 0.20 (5 seeds)   & $+0.173$ (per-seed 95\,\% CIs within $[+0.159,+0.188]$) \\
Genotypes permuted across subjects (5 seeds)  & $-0.001$ \\
Planted genetic share set to zero (3 seeds)   & $-0.001$ \\
\bottomrule
\end{tabular}
\end{table}

The pipeline finds the signal when it is planted and reports none when the
genotype--face pairing is broken or the genetic share is zero (Table
\ref{tab:rq5-controls}). The permutation control was also run on the deep
models: the Transformer scored $0.316$, the MLP $0.252$ and the contrastive
encoder $0.247$, against $0.317$ for demographics. Two points follow. First, no
model gains from uninformative SNPs, which rules out leakage. Second, the MLP
and contrastive encoder lose about $0.07$ when given uninformative SNPs, so on
data with weak genetic signal they could fall below the demographic baseline.
The Transformer's permuted score equals its unpermuted score, which indicates
that at this sample size it is not using the SNPs at all.

\subsection{Effect of sample size}
\label{sec:rq5-scaling}

To ask whether more subjects would help, we generated a separate dataset of
$N=10{,}000$ with the same settings and trained on nested subsets, keeping one
fixed test set of $1{,}500$ (Table \ref{tab:rq5-scaling}). This is a scaling
experiment only: the real cohort has $3{,}118$ subjects, so results beyond the
first row describe a cohort we will not have.

\begin{table}[t]
\centering
\caption{Learning curve on the $N=10{,}000$ scaling dataset, screened SNPs,
test $R^2$ (mean $\pm$ s.d.\ over three training seeds).}
\label{tab:rq5-scaling}
\begin{tabular}{rrrrrr}
\toprule
Train $N$ & Demographics & MLP & Contrastive & Transformer & Diffusion \\
\midrule
$2{,}182$ & $0.312$ & $0.389\pm0.001$ & $0.349\pm0.002$ & $0.310\pm0.001$ & $0.058\pm0.003$ \\
$4{,}000$ & $0.313$ & $0.417\pm0.001$ & $0.392\pm0.001$ & $0.316\pm0.002$ & $0.049\pm0.012$ \\
$7{,}000$ & $0.314$ & $0.435\pm0.001$ & $0.428\pm0.000$ & $0.338\pm0.003$ & $0.028\pm0.011$ \\
\bottomrule
\end{tabular}
\end{table}

The demographic baseline is flat, as expected for a five-parameter linear model.
The MLP gains $+0.046$ and the contrastive encoder $+0.079$ from roughly three
times the data, and the Transformer begins to exceed the baseline only at
$7{,}000$. The diffusion decoder does not improve; we have not diagnosed why and
report it as an open negative result. A companion power grid with the linear
screen gives the same picture: at $N=3{,}118$ the SNPs add $+0.070$ when the
planted genetic share is 0.10 and $+0.177$ when it is 0.20, and nothing when it
is zero.

\paragraph{Limits.}
The simulation is easier than real data in ways that favour detection: 5{,}000
unlinked SNPs instead of roughly $10^6$ correlated ones, 150 causal SNPs that
each affect every shape dimension, and three cleanly separated ancestry groups.
The power figures are therefore an optimistic upper bound. What this section
establishes is that the pipeline is correct, leak-free and sensitive to planted
signal, and that the model comparison is ready to be repeated unchanged on real
data.
